\pagenumbering{arabic}

\chapter{\sc Constructing Tiling Spaces}

We have different methods of constructing tiling spaces.  The first is to be generated by a set of tilings.  If we take $\reals^n$ and subdivide it into different pieces, then each individual piece is called a $\mathit{tile}$, where each tile overlaps another only at its border.  The union of all the tiles is called a $\mathit{tiling}$. For our examples we will only have a finite number of tile types, up to translation, called prototiles.  Each prototile is a polytope.  For a simple example let $n=1$, and let  $T= \cup_{i\in \ints} \ t_i$, where $t_i = [i, i+1]$.  Then each tile only intersects another tile at its boundary points and $T=\cup \ t_i = \reals$.
 A $\mathit{translation}$ of a tiling T is a shift of the tiling by some vector $x\in \reals^n$.  Let T be as above, then $T-1$ is a translation of T to the right by one unit.  
 
 We can place a metric on tilings by considering two tilings close (within $\epsilon$) if after a small translation ($\epsilon$) these tiles agree on a large radius (1/$\epsilon$) about the origin.  It is difficult to state the metric formally.  In the following definition, the closure and therefore, the topology, is given by this metric.
 
 \begin{definition}
 A tiling space $\Omega_T$ is a set of tilings that has the following two properties
 \\ \indent 1) For each $T \in \Omega_T$ and each $x \in \reals^n$ implies $T-x \in$ $\Omega_T$
 \\ \indent 2) $\overline{ \{ T-x \ | \  x \in \reals^n \} }  \subseteq  \Omega_T$ for each T $\in \Omega_T$
 
 \end{definition}
 

